\documentclass[11pt]{article}
\usepackage[T1]{fontenc}
\usepackage[utf8]{inputenc}
\usepackage{lmodern,amsmath,amssymb,amsthm}
\usepackage[margin=1in]{geometry}
\usepackage[hidelinks]{hyperref}
\newtheorem{theorem}{Theorem}
\newtheorem{lemma}[theorem]{Lemma}
\newcommand{\R}{\mathbb R}
\newcommand{\Z}{\mathbb Z}
\newcommand{\N}{\mathbb N}
\newcommand{\poly}{\operatorname{poly}}
\newcommand{\fg}{\mathfrak g}
\newcommand{\id}{\mathrm{id}}
\title{Same-step linearization using the increment polynomial group}
\author{Project proof draft: /root/astra\_v14\_baseline\_theta}
\date{7 November 2026}
\begin{document}
\maketitle

\noindent\textbf{Status.} This is a proposed local replacement for the
same-step linearization construction in GTZ Appendix C. It is not an
independent acceptance. It uses the bounded normalization and tautification
proved in the separately frozen project unit \cite{Taut}, and the published
algebraic polynomial-map results specified below. It does not assume GI($s$).

\section{The required conclusion}

Let $G/\Gamma$ be a connected simply connected nilpotent Lie group modulo a
lattice, with a rational degree filtration $G_\bullet$ of degree at most
$s\geq1$. Every $G_i$ is connected and rational. Let a scalar polynomial
nilsequence be $F(g(n)\Gamma)$, with $F$ complex-valued, $1$-bounded and
Lipschitz. The polynomial $g$ is adapted to $G_\bullet$ in the right-difference
convention $\partial_u g(n)=g(n+u)g(n)^{-1}$.

\begin{theorem}\label{thm:linear}
For this fixed filtered nilmanifold there is a fixed $s$-step nilmanifold
$\widetilde G/\widetilde\Gamma$ with a fixed smooth metric, and a constant
$C$ depending only on the original fixed geometric data, with the following
property. For every polynomial $g$ and every such $F$, there are
$a\in\widetilde G$ and a $1$-bounded Lipschitz function
$\widetilde F:\widetilde G/\widetilde\Gamma\longrightarrow\mathbb C$ such that
\[
 F(g(n)\Gamma)=\widetilde F(a^n\widetilde\Gamma)
 \quad(n\in\Z),\qquad
 \operatorname{Lip}(\widetilde F)\leq C\operatorname{Lip}(F).
\]
Fixed smooth metric comparisons are included in $C$.
The same construction applies to limit polynomial maps and bounded
Lipschitz limit functions, with fixed nilmanifold and uniform constants.
Consequently a finite-family polynomial inverse theorem of degree $s$
implies the finite-family linear inverse theorem of step $s$, without
changing the input Gowers norm or losing correlation.
\end{theorem}

If $\operatorname{Lip}(F)=0$, the constant test function has a constant
extension and the displayed bound is understood literally. A point
nilmanifold can be used for a constant sequence. Step at most $s$ is the
meaning of $s$-step throughout.

By \cite[Theorem 1]{Taut}, restrict to a fixed taut representation with
$G_0=G_1=G$ and replace the polynomial by $h$, where $h(0)=1$.
The exact identity is
\[
 g(0)=q\gamma,\qquad h(n)=q^{-1}g(n)\gamma^{-1},\qquad
 F_q(x\Gamma)=F(qx\Gamma),\qquad
 F_q(h(n)\Gamma)=F(g(n)\Gamma).
\]
Here $q$ lies in a fixed compact set and
$\operatorname{Lip}(F_q)$ is bounded by a fixed multiple of
$\operatorname{Lip}(F)$. The underlying taut filtered nilmanifold is fixed
independently of $g$. It therefore suffices to prove the theorem for this
normalized $h$ and $F_q$. In what follows all geometric data refer to this
fixed taut representation.

\section{Polynomial groups and their Lie algebras}

For $i\geq0$ put
\[
 P_i=\poly(\Z_\N\longrightarrow G_{\bullet+i}).
\]
The zeroth group of the shifted filtration is $G_i$; it need not be taut.
The group operation is pointwise multiplication.
By \cite[Corollary B.4 and Proposition B.6]{GTZ}, these are groups with
\begin{equation}\label{eq:polyfiltered}
 [P_i,P_j]\subseteq P_{i+j},\qquad P_i\triangleleft P_0.
\end{equation}
Those results concern abstract filtered groups and require no equality of
the first two groups. Their conditions hold because the shifted groups
inherit nesting and the commutator law from $G_\bullet$.

Here is an explicit finite-dimensional description, including the smooth
real shift needed below. Let $\fg_i=\log G_i$ and define
\[
 V_i=\left\{X(t)=\sum_{j=0}^{s-i} X_jt^j:X_j\in\fg_{j+i}\right\}
 \quad(0\leq i\leq s),\qquad V_i=0\quad(i>s).
\]
Use pointwise Lie bracket. Bracketing coefficients gives
\begin{equation}\label{eq:Vbracket}
 [V_i,V_j]\subseteq V_{i+j}.
\end{equation}
For instance, a bracket of a degree-$a$ coefficient in $\fg_{a+i}$ and a
degree-$b$ coefficient in $\fg_{b+j}$ belongs to $\fg_{a+b+i+j}$ and has
polynomial degree $a+b$. Thus $V_i$ is a finite-dimensional nilpotent Lie
algebra for every $i$. Its simply connected group is obtained by the
terminating BCH formula, or equivalently by pointwise exponentiation.

\begin{lemma}\label{lem:Pi}
The pointwise exponential identifies that group with $P_i$. In particular
$P_i$ is connected and simply connected, and its smooth topology can be
described either by the coefficients of $\log p(t)$ or by the Taylor
coefficients in \cite[Lemma B.9]{GTZ}.
\end{lemma}
\begin{proof}
If $i>s$, the zeroth group of the shifted filtration is trivial, so both
$P_i$ and $V_i$ are trivial. Fix $0\leq i\leq s$ and put $m=s-i$.
We use the connectedness of every $G_j$ stated in the setup.

First recall explicitly why the global exponential coordinates used here
are smooth. For a connected nilpotent Lie group $H$, let $Q_a(x)=x^a$,
where $a$ is a positive integer. In left tangent trivializations,
differentiating the $a$ factors gives
\[
 DQ_a(x)=\sum_{k=0}^{a-1}\operatorname{Ad}(x)^{-k}.
\]
The infinitesimal adjoint action lowers the lower central series.
An exponential neighborhood generates $H$, so $\operatorname{Ad}(x)$
induces the identity on each successive quotient of that series.
Consequently $DQ_a(x)$ induces multiplication by $a$ on each quotient
and is invertible. The usual local Lie exponential is smooth with
derivative the identity at zero. Given $X\in\mathfrak h$, choose $a$ so
large that $X/a$ lies in this local chart. The identity
\[
 \exp_H(Y)=Q_a(\exp_H(Y/a))
\]
then proves that the exponential has invertible derivative near $X$.
For simply connected nilpotent $H$, its global bijectivity is the
exponential homeomorphism recalled in \cite[Section 3]{GTZ}.
Its local smooth inverses therefore agree to give a smooth global inverse.
This argument uses the local Lie exponential, differentiation of group
operations and the inverse function theorem. It does not infer
smoothness from the word \emph{homeomorphism}.

By \cite[Proposition B.6]{GTZ}, $P_i$ is a group under pointwise
multiplication. By \cite[Lemma B.9]{GTZ}, its elements have the unique
increasing-order Taylor expansion
\[
 p(n)=\prod_{j=0}^{m}u_j^{\binom nj},\qquad u_j\in G_{i+j},
 \qquad n\in\Z.
\]
The coefficients are recovered by the triangular formulas
\[
 u_0=p(0),\qquad
 u_j=\left(\prod_{k<j}u_k^{\binom jk}\right)^{-1}p(j).
\]
The order of the displayed inverse product is part of the formula.
Transfer the product manifold structure of $\prod_{j=0}^m G_{i+j}$
to $P_i$ using these coordinates. Each evaluation at an integer is
smooth. Evaluation of a product or inverse followed by the recovery
formulas proves smoothness of the group operations. The recovered
coefficients lie in the prescribed embedded subgroups; smoothness into
those subgroups follows from their local embedded charts.
The same formulas identify this topology with the pointwise topology.

Evaluation at $0,\ldots,m$ is an injective group homomorphism
$P_i\longrightarrow G^{m+1}$. Hence $P_i$ is nilpotent, with class at most
the class of $G$. Each closed connected subgroup of a simply connected
nilpotent group is simply connected. In exponential coordinates it is
the exponential of its Lie subalgebra: the finite nilpotent BCH formula
makes that exponential image a connected subgroup, and uniqueness of
the connected subgroup with that Lie algebra gives the assertion.
Thus the product manifold above, and therefore $P_i$, is connected and
simply connected. This step does not assert that $P_0$ has class at most $s$.

At the identity, differentiating the Taylor parametrization gives
\[
 (X_0,\ldots,X_m)\longmapsto
 \left(n\longmapsto\sum_{j=0}^m\binom nj X_j\right),
 \qquad X_j\in\fg_{i+j}.
\]
The binomial-to-monomial matrix is triangular with nonzero diagonal.
A contribution to degree $b$ comes only from $j\geq b$, and
$\fg_{i+j}\subseteq\fg_{i+b}$. Descending triangular inversion preserves
these same coefficient spaces. This identifies the tangent space
with exactly $V_i$.
Evaluation maps are smooth homomorphisms, so their derivatives preserve
Lie brackets. The finitely many evaluations just used determine every
tangent vector. They therefore identify the bracket with the pointwise
bracket on $V_i$.

Uniqueness of one-parameter subgroups now gives
\[
 \operatorname{ev}_n(\exp_{P_i}X)=\exp_G(X(n))\qquad(n\in\Z).
\]
Apply the smooth-global-exponential argument above to $P_i$.
Its inverse assigns to $p$ exactly the Lie polynomial whose values
are $\log_G(p(n))$. Conversely every $X\in V_i$ gives an element of $P_i$.
This proves the pointwise exponential identification and the smooth
logarithmic coefficient coordinates. The group law is determined by
multiplication in $G$ at the integer arguments. Under the ordinary
nilpotent BCH identification it is the pointwise BCH law.
The coordinate spaces and finite evaluation maps depend only on the
fixed filtration and $i$.
\end{proof}

For real $u$, define $T^u$ on $V_i$ by $T^uX(t)=X(t+u)$.
Coefficient expansion and nesting show that this preserves $V_i$.
It is a smooth one-parameter group of Lie algebra automorphisms, since
the bracket is pointwise. Exponentiating defines a smooth automorphism
group on $P_i$, which at integer $u$ is the ordinary shift of polynomial
sequences. Its infinitesimal action is ordinary differentiation, and
\begin{equation}\label{eq:derivative}
 \frac{d}{dt}V_i\subseteq V_{i+1}.
\end{equation}
Indeed the coefficient of $t^j$ in $X'(t)$ belongs to
$\fg_{j+1+i}$.

\section{The lattice and the same-step semidirect product}

Define
\[
 \Lambda_i=\{p\in P_i:p(n)\in\Gamma\ \hbox{for every }n\in\Z\}.
\]
It is a subgroup by pointwise multiplication. We give the needed lattice
argument to avoid relying on the disputed step estimate in Appendix C.

\begin{lemma}\label{lem:lattice}
$\Lambda_i$ is a discrete cocompact subgroup of $P_i$. Integer shifts
preserve $\Lambda_i$.
\end{lemma}
\begin{proof}
It is enough to consider $i\leq s$ and put $d=s-i$; for larger $i$ both
groups are trivial. Discreteness follows from evaluation at
$0,\ldots,d$. A lattice-valued polynomial close enough to the identity
at these finitely many points equals the identity at all those points,
since $\Gamma$ is discrete. Taylor uniqueness then makes it the identity
polynomial.

For cocompactness, use induction to write
\[
 p=\gamma_j h_j c_j,
\]
where $\gamma_j\in\Lambda_i$, $c_j$ ranges over a fixed compact subset
of $P_i$, and $h_j$ is the identity at $0,\ldots,j-1$.
For $j=0$ take $\gamma_0=c_0=1$ and $h_0=p$.
Taylor expansion shows $h_j(j)\in G_{j+i}$.
Choose a fixed compact left fundamental set for the lattice
$G_{j+i}\cap\Gamma$ in $G_{j+i}$ and write
$h_j(j)=\lambda_j v_j$, where $\lambda_j$ is in that lattice and $v_j$
is in the compact set. Define
\[
 \ell_j(t)=\lambda_j^{\binom tj},\qquad
 v_j^*(t)=v_j^{\binom tj}.
\]
Both are in $P_i$, by the same one-parameter difference check as above;
$\ell_j\in\Lambda_i$ because $\binom nj$ is an integer for every
$n\in\Z$. Now put
\[
 \gamma_{j+1}=\gamma_j\ell_j,\qquad
 h_{j+1}=\ell_j^{-1}h_j(v_j^*)^{-1},\qquad
 c_{j+1}=v_j^*c_j.
\]
These formulas preserve the factorization and make $h_{j+1}$ the identity
at $0,\ldots,j$. The possible $v_j^*$ form a compact subset of $P_i$,
so the possible $c_{j+1}$ remain in a fixed compact set.
At $j=d+1$, Taylor uniqueness gives $h_{d+1}=1$.
Thus $p=\gamma c$ with $c$ in a fixed compact set, which proves left
cocompactness; inversion gives right cocompactness. Finally an integer
shift preserves both polynomiality and all integer lattice values, and
its inverse does also.
\end{proof}

Now use the \emph{increment} group $P_1$, not $P_0$, and set
\begin{equation}\label{eq:semi}
 \widetilde G=P_1\rtimes_T\R,
 \qquad (p,u)(q,v)=(p\,T^u q,u+v),
 \qquad \widetilde\Gamma=\Lambda_1\rtimes_T\Z.
\end{equation}
The real action just constructed makes $\widetilde G$ a connected simply
connected Lie group, diffeomorphic to $P_1\times\R$.
Its Lie algebra is $V_1\rtimes\R\tau$, where $[\tau,X]=X'$.
Give this Lie algebra the filtration
\[
 \mathfrak l_0=\mathfrak l_1=V_1\oplus\R\tau,
 \qquad \mathfrak l_j=V_j\quad(j\geq2).
\]
Equations \eqref{eq:Vbracket} and \eqref{eq:derivative} prove
$[\mathfrak l_i,\mathfrak l_j]\subseteq\mathfrak l_{i+j}$ for all indices.
In particular $\mathfrak l_{s+1}=0$. Thus $\widetilde G$ is nilpotent of
step at most $s$. This explicit shifted-filtration calculation is the
same-step bound used in the theorem.

The subgroup $\widetilde\Gamma$ is discrete, because its clock coordinate
is integral and $\Lambda_1$ is discrete. It is cocompact as well.
Given $(p,u)$, left-multiply by $(1,-\lfloor u\rfloor)$, then left-reduce
the resulting polynomial coordinate modulo $\Lambda_1$ using Lemma
\ref{lem:lattice}. The resulting representative lies in the product of a
fixed compact set in $P_1$ and $[0,1]$. This proves left cocompactness;
inversion proves right cocompactness. The operations used for left reduction
belong to $\widetilde\Gamma$ by the integer-shift assertion of the lemma.

\section{The orbit identity and fixed compact fiber}

Let $\theta:\widetilde G/\widetilde\Gamma\longrightarrow\R/\Z$ be the
map induced by the clock coordinate, and set $K=\theta^{-1}(0)$.
It is a fixed compact submanifold canonically identified with
$P_1/\Lambda_1$: every element of the fiber can be written
$(p,0)\widetilde\Gamma$, and two such representatives agree precisely when
their polynomial coordinates differ on the right by an element of
$\Lambda_1$.
Evaluation at zero gives a smooth map
\begin{equation}\label{eq:pi}
 \pi:K\longrightarrow G/\Gamma,\qquad
 \pi((p,0)\widetilde\Gamma)=p(0)\Gamma.
\end{equation}
It is well-defined because $\lambda(0)\in\Gamma$ for
$\lambda\in\Lambda_1$. Since $K$ is compact, $\pi$ is Lipschitz for
fixed smooth metrics.

The normalized polynomial $h\in P_0$ has $h(0)=1$. Define its right
increment by
\[
 b(t)=h(t)^{-1}h(t+1).
\]
The left increment $h(t+1)h(t)^{-1}$ lies in $P_1$ directly from the
iterated-difference definition. Proposition B.6 gives $P_1\triangleleft
P_0$, so conjugating that left increment by $h(t)^{-1}$ shows
$b\in P_1$. This verifies the orientation of the increment actually
used, rather than substituting a left increment into a right-telescoping
product.

Put $a=(b,1)\in\widetilde G$. For every integer $n$ define
$b_n(t)=h(t)^{-1}h(t+n)$. The exact cocycle identity
\[
 b_{n+m}(t)=b_n(t)\,b_m(t+n)
\]
and $b_0=1$ prove, for positive and negative integers alike,
\[
 a^n=(b_n,n).
\]
Modulo $\widetilde\Gamma$, right multiplication by $(1,-n)$ gives
$a^n\widetilde\Gamma=(b_n,0)\widetilde\Gamma\in K$.
Consequently
\begin{equation}\label{eq:orbit}
 \pi(a^n\widetilde\Gamma)=b_n(0)\Gamma=h(n)\Gamma,
 \qquad n\in\Z.
\end{equation}
The last equality uses the normalization $h(0)=1$.

\section{Test functions, uniformity and the inverse-theorem join}

The function $F_q\circ\pi$ is $1$-bounded on $K$ and has a Lipschitz
constant bounded by a fixed multiple $L$ of $\operatorname{Lip}(F)$.
Extend its real and imaginary parts separately by the explicit McShane
formula
\[
 E_u(x)=\inf_{y\in K}\bigl(u(y)+L\,d(x,y)\bigr).
\]
This equals $u$ on $K$ and is $L$-Lipschitz: the triangle inequality
proves the Lipschitz assertion, and the original Lipschitz inequality
proves equality at points of $K$. Clip each extension to $[-1,1]$,
combine them as a complex function, and apply metric projection onto
the closed unit disk. This projection $Q$ is $1$-Lipschitz: the
nearest-point inequalities for the disk give
\[
 \langle x-Qx,Qy-Qx\rangle\leq0,\qquad
 \langle y-Qy,Qx-Qy\rangle\leq0,
\]
and therefore $\|Qx-Qy\|^2\leq\langle Qx-Qy,x-y\rangle
\leq\|Qx-Qy\|\|x-y\|$. It also fixes the original values on $K$.
Thus the result $\widetilde F$ is $1$-bounded, agrees with
$F_q\circ\pi$ on $K$, and is at most $\sqrt2 L$-Lipschitz.
Equation \eqref{eq:orbit} proves the exact equality in Theorem
\ref{thm:linear}.

The nilmanifold, lattice, compact fiber, evaluation map and all metric
constants depend only on the fixed original filtered nilmanifold and $s$.
The possibly unbounded polynomial coefficients occur only in the allowed
translation element $a$ and in the bounded representative $q$ affecting the
test function. They do not alter the geometry or its Lipschitz bound.
For limit data the same algebraic formulas and one fixed bounded
representative section give limit choices of $b$ and $a$. Apply the explicit
extension formula to each representative test function. Its uniform bound
and Lipschitz constant give a bounded Lipschitz limit function on the same
standard nilmanifold. The ordinary orbit identities hold at every member
of the defining sequences and therefore at every limit integer.

Suppose now that a polynomial inverse theorem supplies, for fixed $s$ and
$\delta$, a finite family of rational degree-$s$ filtered nilmanifolds,
a fixed Lipschitz bound, and a positive correlation lower bound against
$1$-bounded test functions. Apply the theorem to each member of that finite
family, taking the maximum of its finitely many fixed metric constants.
There is still a finite family, now of step at most $s$ with linear orbits,
and the correlation is unchanged because the represented sequences are
equal pointwise. No comparison of Gowers norms is involved in this join.
This proves the final assertion of Theorem~\ref{thm:linear}.

The construction also covers a point input group: $P_1$ is trivial and the
clock group yields a circle with a constant test function, or it may be
replaced by a point. For $s=1$ the increment group consists of constant
polynomials in an abelian group and the shift action is trivial, so
$\widetilde G$ is abelian, as the required bound predicts.

\section{Comparison with the failed full-group construction}

The original Appendix C used $P_0\rtimes\R$. For the taut abelian
degree-$s$ filtration on $\R$, the commutator of its clock with
$p(t)=\tfrac12\binom ts$, repeated $s$ times, is the nonzero constant
$1/2$. That is a commutator of length $s+1$, so that group is not
$s$-step. This adverse example is recorded in the separate finding packet.
The present construction uses $P_1$ and its increment $b$, reducing
the polynomial degree available in the clock action by one. Its Lie
filtration above verifies the required step bound for all inputs, not
only this example. It changes the construction, while retaining the
same-step linear conclusion needed by GI($s$).

The finite-family polynomial inverse theorem is still a separate premise.
This unit does not prove its rank reduction, symmetrization, or
ultralimit-to-finite-family extraction. Those dependencies remain explicit
in the theorem assembly.

\begin{thebibliography}{9}
\bibitem{GTZ} B. Green, T. Tao and T. Ziegler,
\emph{An inverse theorem for the Gowers $U^{s+1}[N]$-norm},
Annals of Mathematics \textbf{176} (2012), 1231--1372.
Only the specified algebraic polynomial-map and Lie-group facts are used;
the Appendix C step claim is replaced here.
\bibitem{Taut} Project proof, \emph{Tautification and the shifted
nilcharacter construction}, draft 003, 7 November 2026.
The exact author manifest and independent scope decision are pinned in
the accompanying application record.
\end{thebibliography}
\end{document}
